\documentclass[11pt,a4paper]{article}
\usepackage[utf8]{inputenc}
\usepackage[T1]{fontenc}
\usepackage[portuguese]{babel}
\usepackage{geometry}
\geometry{margin=2.5cm}
\usepackage{listings}
\usepackage{xcolor}
\usepackage{amsmath}
\usepackage{amssymb}
\usepackage{hyperref}
\usepackage{enumitem}
\usepackage{tcolorbox}
\usepackage{tabularx}
\usepackage{booktabs}

\definecolor{codebg}{rgb}{0.95,0.95,0.95}
\definecolor{codegreen}{rgb}{0.1,0.5,0.1}
\definecolor{codepurple}{rgb}{0.5,0.1,0.5}
\definecolor{codegray}{rgb}{0.5,0.5,0.5}

\lstdefinestyle{python}{
    backgroundcolor=\color{codebg},
    commentstyle=\color{codegray}\itshape,
    keywordstyle=\color{codepurple}\bfseries,
    stringstyle=\color{codegreen},
    basicstyle=\ttfamily\small,
    breaklines=true,
    frame=single,
    language=Python,
    showstringspaces=false,
    tabsize=4,
    literate={á}{{\'a}}1 {ã}{{\~a}}1 {é}{{\'e}}1 {í}{{\'i}}1 {ó}{{\'o}}1 {õ}{{\~o}}1 {ú}{{\'u}}1 {ç}{{\c{c}}}1 {Á}{{\'A}}1 {Ã}{{\~A}}1 {É}{{\'E}}1 {Í}{{\'I}}1 {Ó}{{\'O}}1 {Õ}{{\~O}}1 {Ú}{{\'U}}1 {Ç}{{\c{C}}}1 {â}{{\^a}}1 {ê}{{\^e}}1 {ô}{{\^o}}1 {û}{{\^u}}1 {Â}{{\^A}}1 {Ê}{{\^E}}1 {Ô}{{\^O}}1 {Û}{{\^U}}1 {à}{{\`a}}1 {è}{{\`e}}1 {ì}{{\`i}}1 {ò}{{\`o}}1 {ù}{{\`u}}1
}
\lstset{style=python}

\newtcolorbox{objetivos}[1][]{
  colback=green!5!white,
  colframe=green!40!black,
  title=O que você vai aprender nesta página,
  fonttitle=\bfseries,
  #1
}

\title{\textbf{Data Analysis with Python 2024--2025}\\Parte 3: Mais NumPy\\\large Processamento de Imagens}
\author{Tradução e Adaptação: University of Helsinki --- MOOC.fi}
\date{}

\begin{document}
\maketitle

\section*{Nota Importante}
Este material é uma tradução e adaptação baseada no curso \textit{Data Analysis with Python 2024-2025}, oferecido pela \textbf{The University of Helsinki MOOC Center}.

\begin{objetivos}
\begin{itemize}
    \item Aprender como processar imagens em Python.
\end{itemize}
\end{objetivos}

\tableofcontents
\newpage

\section{Processamento de Imagens}

\textbf{Nota sobre arquivos:} Para acompanhar os exemplos desta seção e executar os códigos por conta própria, você precisará baixar ou escolher uma imagem qualquer (como uma foto JPG ou PNG) e salvá-la no mesmo diretório em que estiver rodando seus scripts Python. Recomendamos nomear a imagem como \texttt{painting.png} para corresponder aos exemplos abaixo.

Uma imagem é uma coleção de pixels, que é uma abreviação para elementos de imagem (\textit{picture elements}). Uma imagem em tons de cinza pode ser representada como um array bidimensional, cujo primeiro eixo corresponde à coordenada x da imagem e o segundo eixo corresponde à coordenada y. O array contém, em cada par de coordenadas $(x,y)$, um valor que normalmente é um ponto flutuante entre 0.0 e 1.0, ou um inteiro entre 0 e 255. Isso especifica o tom de cinza. Por exemplo, se o array contém o valor 255 nas coordenadas (0,0), então na imagem o pixel no canto superior esquerdo é branco.

Em imagens coloridas, cada pixel tem três componentes de cor (ou canais) correspondentes ao vermelho, verde e azul (RGB). Esses canais formam o terceiro eixo da imagem, com cada canal armazenando um valor que representa a intensidade da respectiva cor. O valor do canal varia entre 0.0 e 1.0 (ou entre 0 e 255). Ao combinar diferentes valores de intensidade para os componentes vermelho, verde e azul de um pixel, é possível representar até 16,7 milhões de cores únicas.

Como as imagens podem ser representadas como arrays multidimensionais, podemos processá-las facilmente usando as funções do NumPy. Vejamos exemplos disso.

\begin{lstlisting}
import matplotlib.pyplot as plt
import numpy as np

painting = plt.imread("painting.png")
print(painting.shape)
print(f"A imagem consiste em {painting.shape[0] * painting.shape[1]} pixels")
plt.imshow(painting)
plt.show()
\end{lstlisting}
\begin{verbatim}
(368, 640, 3)
A imagem consiste em 235520 pixels
\end{verbatim}

Como a imagem agora é um array do NumPy, podemos executar algumas operações facilmente:
\begin{lstlisting}
plt.imshow(painting[:,::-1]); # espelha a imagem na direcao x
\end{lstlisting}

A seguir, definimos os pixels nas primeiras 30 linhas como brancos:
\begin{lstlisting}
painting2 = painting.copy() # nao estrague a pintura original!
painting2[0:30,:,:] = 1.0    # valor maximo para os tres componentes produz branco
plt.imshow(painting2);
\end{lstlisting}

Para uma operação um pouco mais complicada, podemos criar uma função que retorna uma cópia da imagem de modo que ela escureça (fade para o preto) à medida que nos movemos para a esquerda.
\begin{lstlisting}
def fadex(image):
    height, width = image.shape[:2]
    m = np.linspace(0, 1, width).reshape(1, width, 1)
    result = image * m    # note que contamos com o broadcasting aqui
    return result

modified = fadex(painting)
print(modified.shape)
plt.imshow(modified);
\end{lstlisting}
\begin{verbatim}
(368, 640, 3)
\end{verbatim}

\section{Encontrando clusters em uma imagem}

Vamos primeiro gerar alguns dados para este exemplo:
\begin{lstlisting}
n = 5
l = 256
im = np.zeros((l, l))
np.random.seed(0)
points = np.random.randint(0, l, (2, n**2)) # amostra n*n pixels
im[points[0], points[1]] = 1
plt.imshow(im);
\end{lstlisting}

\begin{lstlisting}
from scipy import ndimage
im2 = ndimage.gaussian_filter(im, sigma=l/(8.*n)) # borra um pouco a imagem
plt.imshow(im2);
\end{lstlisting}

Vamos tentar encontrar clusters na imagem gerada acima:
\begin{lstlisting}
mask = im2 > im2.mean() # mascara os pixels cuja intensidade esta acima da media
label_im, nb_labels = ndimage.label(mask) # componentes conectados formam clusters
print(f"O numero de clusters e {nb_labels}")
plt.imshow(label_im);
\end{lstlisting}
\begin{verbatim}
O numero de clusters e 12
\end{verbatim}

Embora esse método que usamos tenha sido muito simples, ele ainda pode ser usado, por exemplo, para contar automaticamente o número de pássaros ou estrelas em uma imagem. Obviamente, humanos conseguem fazer isso facilmente, mas quando existem centenas ou milhares de imagens, é melhor usar máquinas para realizar esse trabalho mecânico.

Existe um grande número de aplicações de processamento de imagens, das quais listamos apenas algumas aqui:
\begin{itemize}
    \item redução de ruído (\textit{denoising})
    \item correção de desfoque (\textit{deblurring})
    \item segmentação de imagens
    \item extração de características (\textit{feature extraction})
    \item zoom, rotação
    \item filtragem
\end{itemize}

\section{Bibliotecas Adicionais}
Existem várias bibliotecas escritas em Python que permitem o fácil processamento de imagens. Alguns exemplos destas:
\begin{itemize}
    \item pillow
    \item scikit-image
    \item No Scipy, há o subpacote \texttt{ndimage} que também contém rotinas para processamento de imagens.
\end{itemize}

\vspace{2em}
\noindent\textbf{Créditos:} Este conteúdo foi originalmente criado pela \textit{The University of Helsinki MOOC Center} para o curso \textit{Data Analysis with Python}.
\end{document}
